\section{Introduction}
\label{sec:intro}
Human demonstrations are the most scalable source of data for dexterous manipulation, yet nearly every way we collect them today, whether teleoperation~\cite{aloha, wu2024gello, opentelevision}, wearable exoskeletons~\cite{umi, doglove, dexumi}, or kinesthetic teaching~\cite{sharma2025selfsupervised, chen2025dexforce}, forces humans to act through a robot embodiment, limiting dexterity to what the interface allows. Retargeting natural human motion removes this bottleneck. However, existing methods~\cite{dexpilot, anyteleop} treat retargeting as a kinematic problem of mapping fingertip positions or joint angles. Further, retargeting often requires scene reconstruction followed by heavy optimization in simulation, and even recent contact-centric formulations~\cite{contactretarget, contactanchored, hoiretarget} still presuppose an object model or a simulator. For manipulation with contact, reconstructing kinematics alone is not enough: millimeters of embodiment mismatch decide whether a replayed trajectory hovers, slips over, or crushes the object.

%while demonstrated contact forces are lost entirely. (see~\cref{fig:teaser} \todo{demonstrate squeezing in fig?})


% Human demonstrations are the most scalable source of dexterous manipulation data, yet nearly every way we collect them for robot learning--teleoperation~\cite{aloha}, exoskeletons~\cite{umi}, kinesthetic teaching~\cite{dexforce}--forces humans to act through a robot embodiment constraining dexterity to what the interface permits. Retargeting natural human motion is a viable alternative, but retargeting has two components: \emph{where} the end-effector reaches, and \emph{how} it affects the environment through contact. Existing work~\todo{add citations} casts retargeting as a kinematic problem of fingertip positions or joint angles~\citep{dexpilot, anyteleop}. At contact, this is not enough: millimeters of embodiment error decide whether a replayed grasp hovers, slips, or crushes, and the forces the human applied are never represented at all. 

We present \ours, a retargeting framework that treats \fca as a first-class citizen. Instead of relying on object geometry, \ours uses per-link \emph{contact force targets} extracted from a human demonstration, comprising a center of pressure, a push direction, and a force, to solve for robot configurations that place semantically matched robot links on every contact force target inside the demonstrated friction cone. \ours produces contact force targets alongside the joint trajectory without relying on object information, therefore it needs no simulator, object mesh, or object pose estimate. Consequently, rigid, articulated, and deformable objects need no special handling.

% \ours makes contact a first-class citizen of retargeting. From any demonstration it extracts per-phalanx \emph{contact anchors}---a center of pressure on the hand surface, a push direction, and a force---and solves for robot poses that place the matching fingerpad on each anchor inside the friction cone of the demonstrated push, producing force targets alongside kinematic ones.


\ours is agnostic to the demonstration source: any modality containing hand pose and contact can be used. To make heterogeneous sources interchangeable, we introduce a \emph{unified hand--contact representation} that turns monocular video (through Vid-HACO, our video contact model), tactile gloves, and marker-based motion capture into a common interface: a MANO mesh sequence with one \emph{contact force target} per hand region. Because every contact force target is bound to the human hand surface, retargeting to a new robot requires only a semantic map from hand regions to robot bodies.

On hardware, \ours pairs retargeting with a \emph{tactile tracking controller}, an integral admittance loop over joint impedance that maintains the demonstrated contact forces during replay. Together they enable \emph{open-loop replay} of dexterous human demonstrations on multiple robot embodiments mounted on a 7-DoF arm. Here \emph{open-loop} means no object pose estimate and no learned policy. The only feedback is the fingertip force loop that maintains the demonstrated contact forces. We prioritize open-loop replay to demonstrate the retargeting quality of \ours over policy learning, since a closed-loop policy confounds retargeting fidelity with other issues such as learning error, partial observability and the like. %We replay rather than learn because open-loop replay is the strictest test of retargeting fidelity: no policy can compensate for a wrong contact, and a learned policy would confound retargeting quality with learning.

% Anchoring retargeting to contact also makes it \emph{general across data sources}. Monocular video, tactile gloves, and marker-based motion capture differ wildly as sensors, but each can produce a hand mesh with per-vertex contact; one shared extraction chain turns all of them into the same anchors, and one solver configuration serves all three. For RGB-only data, our video contact model Vid-HACO\todo{confirm Vid-HACO expansion} predicts per-vertex contact directly from monocular video, so ordinary phone recordings of bare-handed manipulation become retargetable demonstrations.

% These sources meet in a \emph{unified hand--contact representation}: a MANO mesh sequence with per-vertex contact probability, from which semantic contact regions and per-link centers of pressure are derived. Because every target lives on the human hand, the representation is embodiment-agnostic---a new robot needs only a semantic map from hand regions to robot bodies. Retargeting then emits \emph{tactile targets} along with the joint trajectory, and a tactile tracking controller---integral admittance on measured fingertip force over joint impedance---holds those targets during rollout, preserving contact where pure position replay would hover or crush. Together they replay highly dexterous demonstrations \emph{open-loop} on a 22-DoF tactile hand mounted on a 7-DoF arm, with no object pose estimation and no policy.

% Finally, the formulation is \emph{scalable}. It needs no simulator and no object model: the object is never reconstructed, posed, or meshed, so interactions with rigid, articulated, and deformable objects---which break object-centric and simulation-based pipelines---require nothing special. A GPU sampling solver retargets at interactive rates, fast enough to also drive live glove teleoperation.

Concretely, our contributions are:
\begin{itemize}
    \item \textbf{Contact-Anchored Retargeting:} A simulator and object-free formulation that optimizes contact reach, friction-cone alignment, force weighting, and contact stability per frame with a fast GPU sampling optimizer.
    \item \textbf{Unified Hand--Contact Representation:} A shared interface that maps video (through Vid-HACO, our video hand-contact model), tactile gloves, and motion capture into contact force targets with whole-clip force synthesis for sources that measure no force.
    \item \textbf{Tactile Tracking Controller:} A real-world force tracking controller with feedforward force that replays contact-rich tasks open-loop faithfully on physical hardware.
    \item \textbf{Cross-Embodiment Generalization:} Cross-embodiment retargeting through a semantic map, shown on four robot hands from a 22-DoF dexterous hand down to a parallel-jaw gripper.
\end{itemize}

Across video, glove, and motion-capture sources, \ours achieves state-of-the-art contact fidelity, matching demonstrated contact on 43--78\% of contact frames where kinematic retargeting reaches under 10\%, while retargeting methods needing object information~\cite{regrind} only reach 22.5\%. Further, our trajectories transfer to the real world through open-loop replay, where the tactile tracker raises success from 4 to 25 of 30 trials over position-only replay.

% Concretely, our contributions are:
% \begin{itemize}[leftmargin=*, itemsep=1pt, topsep=2pt]
%     \item A \textbf{contact-anchored retargeting formulation}---contact reach, friction-cone orientation, force weighting, and contact stability---solved per frame by a GPU sampling optimizer with a deterministic bone-direction initializer, entirely free of simulators and object models.
%     \item A \textbf{unified hand--contact representation} (UHCR) and shared extraction chain that converts monocular video (via our Vid-HACO contact model), tactile-glove, and motion-capture demonstrations into the same contact anchors, including force synthesis with a whole-clip grasp-balance constraint for sources without measured forces.
%     \item A \textbf{real-world tactile tracking controller} that tracks the retargeted tactile targets---contact state machine, integral admittance on fingertip force, and force feedforward---enabling open-loop replay of contact-rich demonstrations on hardware.
%     \item \textbf{Cross-embodiment transfer} through per-robot semantic maps alone, instantiated on four embodiments from a 22-DoF hand down to a parallel-jaw gripper.
% \end{itemize}

% Across video, glove, and motion-capture sources, \ours \todo{headline retargeting result vs.\ kinematic baselines}, and replays \todo{headline hardware result: tasks/success} open-loop on hardware.